\input _template

\setlanguage{CS}

\Title{Algebra}

\Subtitle{Vtipná algebra}

\MathcompsLink{vtipna-algebra}

\Author{Patrik Bak}

\sec Úvod

V této sbírce úloh najdete mé nejoblíbenější úlohy, které mají krátká zajímavá, častokrát triková řešení využívající různé algebraické triky.

\sec Úlohy

Bez dalšího otálení pojďme na úlohy. Klíčem je najít vtipnou úpravu. Je těžké dát obecný návod -- koneckonců, nadstandardní úlohy jsou přece jen dost o kreativitě.

\Problem{0}{}{
    Označme
    $$
    m \circ n = \frac{mn+4}{m+n}.
    $$
    Určete
    $$
    (\dots((42 \circ 41) \circ 40) \circ \dots \circ 1) \circ 0.
    $$
}{
    V~úloze jsou zdánlivě náhodné konstanty, například~$42$ vypadá, že by mohla být uměle nastavená. Kdybychom měli ručně vyhodnocovat takto velká čísla, zblázníme se. Zkuste si úlohu pro menší konstanty.
}{
    Platí
    \begitems
    \i $1 \circ 0 = 4$
    \i $(2 \circ 1) \circ 0 = 2$
    \i $((3 \circ 2) \circ 1) \circ 0 = 2$
    \i $(((4 \circ 3) \circ 2) \circ 1) \circ 0 = 2$
    \enditems
    Není podezřelé, že to po čase vychází stále 2?
}{
    Všimněme si speciální vlastnosti čísla $2$ vzhledem k operaci $\circ$. Pro libovolné $m \ne -2$ platí
    $$
    m \circ 2 = \frac{2m+4}{m+2} = \frac{2(m+2)}{m+2} = 2.
    $$
    Výraz ze zadání vyhodnocujeme postupně zleva, přičemž čísla vstupují do operace v~pořadí $42, 41, \dots, 2, 1, 0$. V~momentě, kdy budeme mezivýsledek skládat s~číslem $2$, se hodnota celého výrazu změní na $2$. Všechny následující operace budou mít tvar $2 \circ k$, což je opět $2$. Výsledek je proto $2$.
}

\Problem{0}{}{
    Čísla $a,b,c$ jsou nenulová reálná čísla se součtem $0$. Dokažte, že
    $$
    \frac{1}{a+b}+\frac{1}{b+c}+\frac{1}{c+a}=\frac{a^2+b^2+c^2}{2abc}.
    $$
}{
    Podmínku $a+b+c=0$ chceme šikovně využít, abychom se zbavili komplikovaných zlomků nalevo.
}{
    Klíčem je uvědomit si, že $a+b=-c$, tedy první zlomek je vlastně roven $-\frac 1c$, podobně zbylé.
}{
    Po úpravě stačí ukázat
    $$
    -\left(\frac 1a + \frac 1b + \frac 1c\right) = \frac{a^2+b^2+c^2}{2abc},
    $$
    po roznásobení
    $$
    2ab+2bc+2ca=a^2+b^2+c^2.
    $$
    Teď potřebujeme znovu použít podmínku $a+b+c=0$. Co takového bychom s~ní mohli udělat, abychom dostali výrazy, o~kterých máme něco dokázat?
}{
    Z~podmínky $a+b+c=0$ vyplývá $a+b=-c$, $b+c=-a$ a~$c+a=-b$. Levou stranu dokazované rovnosti proto upravíme na
    $$
    -\frac{1}{c} - \frac{1}{a} - \frac{1}{b} = -\frac{ab+bc+ca}{abc}.
    $$
    Zároveň umocněním rovnosti $a+b+c=0$ na druhou dostaneme
    $$
    0 = (a+b+c)^2 = a^2+b^2+c^2 + 2(ab+bc+ca),
    $$
    z~čehož vyjádříme $ab+bc+ca = -\frac{1}{2}(a^2+b^2+c^2)$. Dosazením tohoto výrazu do upravené levé strany ihned získáme
    $$
    -\frac{-\frac{1}{2}(a^2+b^2+c^2)}{abc} = \frac{a^2+b^2+c^2}{2abc},
    $$
    což jsme chtěli dokázat.
}

\Problem{0}{}{
    Čísla $a,b,c$ jsou kladná reálná taková, že $ab+bc+ca=1$. Určete všechny možné hodnoty výrazu
    $$
    \frac{a(b^2+1)}{a+b}+
    \frac{b(c^2+1)}{b+c}+
    \frac{c(a^2+1)}{c+a}.
    $$
}{
    Trikem je zjednodušit každý zlomek individuálně použitím podmínky $ab+bc+ca=1$.
}{
    Náš zkoumaný výraz je plný písmen a~je tam jedno číslo, $1$. Co takhle se ho zbavit?
}{
    Upravme první zlomek dosazením $1 = ab+bc+ca$ do čitatele:
    $$
    \gather
    \frac{a(b^2+ab+bc+ca)}{a+b} = \frac{a(b(a+b)+c(a+b))}{a+b} =\\= \frac{a(a+b)(b+c)}{a+b} = ab+ac.
    \endgather
    $$
    Cyklicky pro druhý a~třetí zlomek dostaneme $bc+ab$ a~$ca+bc$. Sečtením všech tří upravených výrazů získáme
    $$
    (ab+ac) + (bc+ab) + (ca+bc) = 2(ab+bc+ca).
    $$
    Jelikož $ab+bc+ca=1$, hodnota výrazu je $2$.

    Ještě potřebujeme dokázat, že tato hodnota je dosažitelná -- zřejmě však existují kladná čísla $a,b,c$, pro která $ab+bc+ca=1$, například $a=b=c=\sqrt{3}/3$.
}

\Problem{0}{}{
    Čísla $a,b,c$ jsou kladná reálná se součinem 1. Určete všechny možné hodnoty výrazu
    $$
    \frac{1}{1+a+ab}+
    \frac{1}{1+b+bc}+
    \frac{1}{1+c+ca}.
    $$
}{
    Klíčem je provést nějakou šikovnou úpravu některých zlomků, která umožní použít podmínku $abc=1$. Co je taková běžná věc, kterou umíme dělat se zlomky?
}{
    Jednoduchý způsob, jak využít $abc=1$ např. v~prvním zlomku, je rozšířit ho~$c$, potom máme
    $$
    \frac{1}{c+ac+abc}=\frac{1}{c+ac+1}.
    $$
    Tento zlomek má čistou náhodou stejný jmenovatel jako poslední zlomek.
}{
    První zlomek rozšíříme číslem $c$ a~druhý výrazem $ac$. S~využitím $abc=1$ tak upravíme jejich jmenovatele na tvar shodný se třetím zlomkem:
    $$
    \align
    \frac{1}{1+a+ab} &= \frac{c}{c+ac+abc} = \frac{c}{c+ac+1}, \\
    \frac{1}{1+b+bc} &= \frac{ac}{ac+abc+ab c^2} = \frac{ac}{ac+1+c}.
    \endalign
    $$
    Sečtením všech tří zlomků dostaneme
    $$
    \frac{c}{1+c+ac} + \frac{ac}{1+c+ac} + \frac{1}{1+c+ac} = \frac{c+ac+1}{1+c+ac} = 1.
    $$
    Jediná možná hodnota výrazu je $1$. Příklad $a=b=c=1$ přitom dokazuje, že tato hodnota je dosažitelná.
}

\Problem{0}{}{
    Reálná čísla $a, b, c$ mají součet 1. Dokažte, že číslo
    $$
    (a + bc)(b + ca)(c + ab)
    $$
    je nezáporné.
}{
    Trikem je upravit každou závorku individuálně použitím podmínky $a+b+c=1$.
}{
    Dobrým indikátorem, jak upravit každou závorku, je fakt, že v~jednotlivých závorkách máme $a+bc$, kde je první člen stupně~1 a~druhý stupně~2 -- co takhle je nějak oba udělat stupně~2?
}{
    Využijeme podmínku $a+b+c=1$ na úpravu výrazů v~závorce. Pro první závorku platí:
    $$
    \gather
    a + bc = a(a+b+c) + bc = a^2+ab+ac+bc =\\= a(a+b)+c(a+b) = (a+c)(a+b).
    \endgather
    $$
    Analogicky odvodíme $b+ca = (a+b)(b+c)$ a~$c+ab = (a+c)(b+c)$. Vynásobením všech tří výrazů dostaneme
    $$
    (a+bc)(b+ca)(c+ab) = (a+b)^2(b+c)^2(c+a)^2.
    $$
    Výsledek je součinem druhých mocnin reálných čísel, a~proto je nezáporný.
}

\Problem{0}{}{
    Pro nenulová reálná čísla $a,b,c$ platí
    $$
    \frac1a + \frac1b + \frac1c = \frac1{a+b+c}.
    $$
    Dokažte, že pro každé liché přirozené číslo $n$ platí
    $$
    \frac1{a^n} + \frac1{b^n} + \frac1{c^n} = \frac1{a^n+b^n+c^n}.
    $$
}{
    Zapomeňte na to, co dokazujeme. Kdy vůbec může platit předpoklad? Zkuste najít všechny trojice $a,b,c$, které tuto rovnost splňují.
}{
    Klíčem je uvědomit si, že po roznásobení podmínky dostaneme
    $$
    (a+b+c)(ab+bc+ca)-abc=0.
    $$
    Tento výraz se překvapivě dá rozložit na součin. Odtud to už bude přímočaré.
}{
    Upravme podmínku ze zadání vynásobením jmenovateli:
    $$
    \gather
    \frac{ab+bc+ca}{abc} = \frac{1}{a+b+c},\\
    (a+b+c)(ab+bc+ca) - abc = 0.
    \endgather
    $$
    Výraz na levé straně upravíme postupným vytýkáním na součin:
    $$
    \align
    (a+b+c)(ab+bc+ca) - abc &= \\ =
    (a+b+c)(ab + c(a+b)) - abc &= \\ =
    (a+b)ab + abc + (a+b)^2c + (a+b)c^2 - abc &= \\ =
    (a+b)(ab + ac + bc + c^2) &= \\ =
    (a+b)(a(b+c) + c(b+c))  &= \\ =
    (a+b)(b+c)(c+a).&
    \endalign
    $$
    Rovnost tedy platí právě tehdy, když jsou alespoň dvě z~čísel opačná. Bez újmy na obecnosti nechť $b=-a$.
    Potom pro každé liché přirozené číslo $n$ platí $b^n = (-a)^n = -a^n$.
    Levá strana dokazované rovnosti je
    $$
    \frac{1}{a^n} + \frac{1}{-a^n} + \frac{1}{c^n} = \frac{1}{c^n}.
    $$
    Pravá strana je
    $$
    \frac{1}{a^n - a^n + c^n} = \frac{1}{c^n}.
    $$
    Levá strana se rovná pravé, čímž je důkaz hotov.
}

\Problem{0}{}{
    Navzájem různá reálná čísla $a,b,c$ splňují
    $$
    a + \frac1b = b + \frac1c = c + \frac1a.
    $$
    Dokažte, že $|abc|=1$.
}{
    Zadání nám říká, že platí
    $$
    \align 
    a + \frac 1b &= b + \frac 1c, \\
    b + \frac 1c &= c + \frac 1a, \\
    c + \frac 1a &= a + \frac 1b. \\
    \endalign
    $$
    Uvědomění si, že jde o~soustavu tří rovnic, je dobrý krok k~tomu, aby bylo přirozené zkoušet různě kombinovat rovnice, jak je běžné při práci se soustavami.
}{
    Po chvíli snadno zjistíme, že sčítání a~odčítání rovnic k~ničemu nevede. Klíčem je tyto rovnice násobit. K~tomu však potřebujeme každou rovnici trochu upravit, aby se nám při násobení něco vykrátilo.
}{
    Z~rovnosti $a + \frac{1}{b} = b + \frac{1}{c}$ vyjádříme rozdíl $a-b$:
    $$
    a - b = \frac{1}{c} - \frac{1}{b} = \frac{b-c}{bc}.
    $$
    Cyklicky pro další dvojice získáme
    $$
    b - c = \frac{c-a}{ca} \quad\text{a}\quad c - a = \frac{a-b}{ab}.
    $$
    Jelikož jsou čísla $a,b,c$ navzájem různá, rozdíly $a-b$, $b-c$ a~$c-a$ jsou nenulové. Vynásobením těchto tří rovností dostaneme
    $$
    (a-b)(b-c)(c-a) = \frac{(b-c)(c-a)(a-b)}{(abc)^2}.
    $$
    Vydělením nenulovým výrazem $(a-b)(b-c)(c-a)$ získáme $1 = \frac{1}{(abc)^2}$, tedy $(abc)^2 = 1$, z~čehož vyplývá $|abc|=1$.
}

\Problem{0}{}{
    Pro nenulová reálná čísla $a,b,c$ platí
    $$
    a^2-b^2 = bc \hbox{\quad a \quad} b^2-c^2 = ca.
    $$
    Dokažte, že $a^2 - c^2 = ab$.
}{
    Sečtením rovnic dostaneme přesně levou stranu dokazované rovnosti. Stačí tedy dokázat něco, co vypadá zdánlivě jednodušeji.
}{
    Podle prvního návodu jsme už zjistili, že součet rovnic z~podmínky má smysl. To ještě není všechno, potřebujeme něco dalšího. Jiné než sčítání rovnic je například násobení. Předtím je ale často musíme upravit -- ideálně tak, aby se po násobení něco vykrátilo.
}{
    Sečtením zadaných rovnic dostaneme
    $$
    a^2-c^2 = bc+ca.
    $$
    K důkazu tvrzení tedy stačí ukázat, že $ab = bc+ca$.
    První rovnici ze zadání upravíme na $a^2 = b(b+c)$ a~druhou na $(b-c)(b+c)=ca$. Vynásobením těchto rovností máme
    $$
    a^2(b-c)(b+c) = abc(b+c).
    $$
    Kdyby $b+c=0$, tak z~první upravené rovnice $a^2=0$, tedy $a=0$, což odporuje zadání. Proto můžeme rovnici vydělit nenulovým výrazem $a(b+c)$, čímž dostaneme
    $$
    a(b-c) = bc,
    $$
    což je zřejmě ekvivalentní dokazovanému vztahu $ab=bc+ca$. Tím je důkaz hotov.
}

\Problem{1}{}{
    Jen s~použitím pera a~papíru (a mozku) najděte dvě čtyřmístná čísla, jejichž součin je roven $4^8 + 6^8 + 9^8$.
}{
    Úloha nám naznačuje, že máme hledat rozklad na součin. Náš výraz je ale velmi nepřehledný, zaveďme nejprve vhodnou substituci čísel za písmenka, abychom ho zpřehlednili.
}{
    Vyjádření $4=2^2$, $9=3^2$ a $6=2 \cdot 3$ naznačují, že klíčová substituce je $a=2$, $b=3$, výraz je potom roven
    $$
    a^{16}+(ab)^8+b^{16}.
    $$
    Dokážeme toto rozložit na součin?
}{
    Trikem k~rozkladu $a^{16}+(ab)^8+b^{16}$ je \textit{doplnění na čtverec}, zkuste doplnit $a^{16}+b^{16}$ na čtverec.
}{
    Nechť $a=2$, $b=3$. Výraz ze zadání je potom roven
    $$
    a^{16} + a^8 b^8 + b^{16}.
    $$
    Tento výraz upravíme doplněním na čtverec a~následným použitím vzorce pro rozdíl čtverců:
    $$
    \align
    a^{16}+a^8b^8+b^{16} &= (a^{16} + 2a^8b^8 + b^{16}) - a^8b^8 \\
    &= (a^8+b^8)^2 - (a^4b^4)^2 \\
    &= (a^8+b^8-a^4b^4)(a^8+b^8+a^4b^4).
    \endalign
    $$
    Dosadíme $a=2$ a~$b=3$. Potom $a^4=16$, $b^4=81$, $a^4b^4 = 1296$.
    Dále $a^8=256$ a~$b^8=6561$, a $a^8+b^8 = 256+6561 = 6817$.
    Jednotlivé závorky jsou tedy rovny:
    $$
    \align
    6817 - 1296 &= 5521, \\
    6817 + 1296 &= 8113.
    \endalign
    $$
    Obě nalezená čísla jsou čtyřmístná.
}


\Problem{1}{}{
    V~oboru reálných čísel řešte rovnici
    $$
    x^4+4x^3+6x^2+4x=6.
    $$
}{
    Konstanty na levé straně nejsou náhodné a~měly by něco připomínat.
}{
    To něco, co by tyto konstanty měly připomínat, se dokonce učí ve škole a má jméno.
}{
    Levá strana rovnice připomíná binomickou větu pro 4~členy. Vskutku,
    $$
    (x+1)^4 = x^4+4x^3+6x^2+4x+1.
    $$
    To je skoro levá strana rovnice, potřebujeme ještě přičíst~1. Rovnice je tedy ekvivalentní s~rovnicí.
    $$
    (x+1)^4 = 7.
    $$
    V~oboru reálných čísel má rovnice dvě řešení $x=-1\pm \root 4 \of 7$.
}

\Problem{2}{}{
    V~oboru reálných čísel řešte rovnici
    $$
    (x^2 + 3x + 2)(x^2 - 2x - 1)(x^2 - 7x + 12) + 24 = 0.
    $$
}{
    Kvadratické členy nejsou náhodné. Co takhle se podívat na každý zvlášť, nedá se přetransformovat?
}{
    Prvním krokem úlohy je uvědomit si, že dva ze tří našich kvadratických členů se dají rozložit na součin -- kupodivu to vychází i pěkně, je to náhoda?
}{
    Co se čtyřmi lineárními členy a~jedním kvadratickým? Trikem je z~nich znovu vytvořit tři kvadratické členy.
}{
    Aplikováním rad z~předešlých návodů bychom měli dojít k~rovnici
    $$
    (x^2 - 2x - 8)(x^2 - 2x - 3)(x^2 - 2x - 1) + 24=0.
    $$
    Její stupeň umíme zjednodušit vhodnou substitucí. Měla by zůstat rovnice s~celočíselnými kořeny, kterou už vyřešíme středoškolskými postupy.
}{
    Rozložíme první a~třetí závorku na součin lineárních činitelů:
    $$
    (x+1)(x+2)(x^2-2x-1)(x-3)(x-4) + 24 = 0.
    $$
    Vhodným přeskupením členů máme:
    $$
    \align
    (x+1)(x-3) &= x^2-2x-3 \\
    (x+2)(x-4) &= x^2-2x-8
    \endalign
    $$
    Dosazením do původní rovnice máme
    $$
    (x^2-2x-3)(x^2-2x-8)(x^2-2x-1) + 24 = 0.
    $$
    Zavedeme substituci $y = x^2-2x$. Rovnice přejde do tvaru
    $$
    (y-3)(y-8)(y-1) + 24 = 0.
    $$
    Po roznásobení a~úpravě získáme
    $$
    \align
    y^3 - 12y^2 + 35y &= 0 \\
    y(y^2-12y+35) &= 0 \\
    y(y-5)(y-7) &= 0.
    \endalign
    $$
    Pro $y \in \{0, 5, 7\}$ dořešíme kvadratické rovnice pro $x$:
    \begitems
    \i $x^2-2x = 0$ má kořeny $0, 2$.
    \i $x^2-2x-5 = 0$ má kořeny $1 \pm \sqrt{6}$.
    \i $x^2-2x-7 = 0$ má kořeny $1 \pm 2\sqrt{2}$.
    \enditems
    Množina řešení je $\{0, 2, 1 \pm \sqrt{6}, 1 \pm 2\sqrt{2}\}$.
}


\Problem{2}{}{
    Celá čísla $a,b,c$ splňují
    $$
    (a-b)^2 + (b-c)^2 + (c-a)^2 = abc.
    $$
    Dokažte, že $a+b+c+6 \mid a^3+b^3+c^3$.
}{
    Klíčem je vzpomenout si na algebraickou identitu, která jaksi spojuje některé z~výrazů, o~kterých je tato úloha.
}{
    Klíčová identita je rozklad výrazu $a^3+b^3+c^3-3abc$ na součin. Pokud to neznáte, zkuste si tento rozklad odvodit. V~dalším návodu prozradíme možné triky, jak na něj přijít.
}{
    K~objevení $a^3+b^3+c^3-3abc$ můžeme použít tyto možné postupy:
    \begitems \style a
    \i Doplníme $a^3+b^3$ na třetí mocninu, tedy vyjádříme z~rozkladu $(a+b)^3$.
    \i Doplníme $a^3+b^3+c^3$ na třetí mocninu, tedy vyjádříme z~rozkladu $(a+b+c)^3=...$. Z~toho extrahujeme $a^3+b^3+c^3=...$. Zůstanou nám pouze smíšené členy jako $a^2b$, $ab^2$ a $abc$. Ty je potřeba vhodně uspořádat, abychom vytkli něco před závorku.
    \enditems 
}{
    Prozradíme klíčovou identitu
    $$
    a^3+b^3+c^3-3abc = (a+b+c)(a^2+b^2+c^2-ab-bc-ca).
    $$
    (V~předešlém návodu naznačené odvození najdete v~řešení.)
}{
    Jak využít klíčovou identitu? Potřebujeme ještě jeden trik: uvědomění si, jak druhá závorka rozkladu souvisí s~podmínkou ze zadání.
}{
    Podmínku ze zadání můžeme po roznásobení napsat jako
    $$
    2(a^2+b^2+c^2-ab-bc-ca)=abc.
    $$
    Shodou okolností se nám tu objevila druhá závorka našeho rozkladu. Teď už stačí dát všechno opatrně dohromady.
}{
    Nejprve odvodíme identitu pro rozklad výrazu $a^3+b^3+c^3-3abc$. Součet prvních dvou třetích mocnin doplníme na třetí mocninu součtu podle vztahu
    $$
    a^3+b^3 = (a+b)^3 - 3ab(a+b).
    $$
    Následně použijeme vzorec pro součet třetích mocnin
    $$
    x^3+y^3=(x+y)(x^2-xy+y^2)
    $$
    na členy $(a+b)^3$ a~$c^3$:
    $$
    \align
    a^3+b^3+c^3-3abc &=\\ 
    (a+b)^3 + c^3 - 3ab(a+b) - 3abc &=\\
    (a+b+c)((a+b)^2 - (a+b)c + c^2) - 3ab(a+b+c) &=\\
    (a+b+c)(a^2+2ab+b^2 - ac - bc + c^2 - 3ab) &= \\
    (a+b+c)(a^2+b^2+c^2-ab-bc-ca)&. 
    \endalign
    $$
    Druhou závorku označme $K = a^2+b^2+c^2-ab-bc-ca$. Jelikož $a,b,c$ jsou celá čísla, i $K$ je celé číslo.
    Podmínku ze zadání roznásobíme a~upravíme do tvaru
    $$
    (a^2-2ab+b^2) + (b^2-2bc+c^2) + (c^2-2ca+a^2) = abc,
    $$
    takže $2K=abc$. Tento vztah dosadíme do odvozené identity:
    $$
    a^3+b^3+c^3 - 3(2K) = (a+b+c)K.
    $$
    Osamostatněním součtu třetích mocnin dostaneme
    $$
    a^3+b^3+c^3 = 6K + (a+b+c)K = K(a+b+c+6).
    $$
    Výraz $a+b+c+6$ tedy dělí $a^3+b^3+c^3$.

    \textit{Poznámka.} Ukážeme ještě alternativní způsob doplnění na třetí mocninu pomocí rozkladu ${(a+b+c)^3}$. Po roznásobení platí:
    $$
    \gather
    (a+b+c)^3 = a^3+b^3+c^3+ \\
    {}+ 3(a^2b+ab^2+b^2c+bc^2+c^2a+ca^2) + 6abc
    \endgather
    $$
    takže
    $$
    \gather
    a^3+b^3+c^3 - 3abc = (a+b+c)^3 \\
    {}- 3(a^2b+ab^2+b^2c+bc^2+c^2a+ca^2)- 9abc,
    \endgather
    $$
    Z~posledního řádku lze vytknout~$-3$ před závorku, soustřeďme se pouze na vnitřek této závorky:
    $$
    \align
    (a^2b+ab^2+b^2c+bc^2+c^2a+ac^2) + 3abc &= \\
    a^2b+ab^2+b^2c+c^2b+c^2a+a^2c+3abc &= \\
    [a^2b+ab^2+abc]+[b^2c+c^2b+abc]+[c^2a+a^2c+abc] &=\\
    ab(a+b+c) + bc(b+c+a) + ca(c+a+b) &= \\
    (a+b+c)(ab+bc+ca).
    \endalign
    $$
    Teď to už jen dát dohromady:
    $$
    \gather
    a^3+b^3+c^3 - 3abc = (a+b+c)^3 \\
    {}- 3(a^2b+ab^2+b^2c+bc^2+c^2a+ca^2)- 9abc = \\
    (a+b+c)^3 - 3(a+b+c)(ab+bc+ca) \\ 
    (a+b+c)((a+b+c)^2-3(ab+bc+ca)),
    \endgather
    $$
    což je alternativní zápis naší dokazované identity.
}

\DisplayHints

\DisplayProblemSolutions

\bye